\documentclass[12pt,a4paper]{amsart}
\usepackage{amsmath,amssymb,amsthm}
\usepackage{mathtools}
\usepackage[utf8]{inputenc}
\usepackage[T1]{fontenc}
\usepackage{hyperref}
\usepackage{enumitem,booktabs,array}
\usepackage{geometry}
\geometry{margin=2.5cm}

\newtheorem{theorem}{Theorem}[section]
\newtheorem{lemma}[theorem]{Lemma}
\newtheorem{corollary}[theorem]{Corollary}
\newtheorem{proposition}[theorem]{Proposition}
\newtheorem{conjecture}[theorem]{Conjecture}
\theoremstyle{definition}
\newtheorem{definition}[theorem]{Definition}
\newtheorem{example}[theorem]{Example}
\newtheorem{remark}[theorem]{Remark}

\author{Michael Fuhrmann}
\address{Specialist Mathematics Project, Build 138}
\email{specialist-maths@localhost}
\date{\today}

\title{Differential Geometry: Manifolds, Curvature, and the Interplay of\\
Geometry and Analysis}

\begin{abstract}
Differential geometry is the mathematical discipline that studies smooth
manifolds and their geometric properties using the tools of differential
calculus and tensor analysis. This paper provides a rigorous yet accessible
introduction to the core concepts of modern differential geometry: smooth
manifolds with charts and atlases, Riemannian metrics defining lengths and
angles, connections and covariant derivatives leading to the fundamental
Levi-Civita connection (proved unique via the metric compatibility and
torsion-free conditions), and the Riemann curvature tensor with its derived
quantities --- the Ricci tensor, scalar curvature, and Gaussian curvature.
The classical Gauss-Bonnet theorem, which links geometry to topology via
the Euler characteristic, is stated and proved for closed surfaces; its
far-reaching generalisation, the Chern-Gauss-Bonnet theorem for even-dimensional
manifolds, is also presented. We treat geodesics through the geodesic
equation and the Hopf-Rinow theorem (proved), Jacobi fields, and the
second fundamental form for submanifolds. Lie groups as geometric objects
equipped with left-invariant metrics provide an algebraic perspective.
The paper concludes with physical applications in general relativity, where
the Einstein field equations describe spacetime curvature sourced by
matter and energy. All results are illustrated by concrete examples:
spheres, tori, hyperbolic planes, and minimal surfaces. The accompanying
Python module \texttt{differential\_geometry.py} implements the algorithms
discussed here, providing numerical verification of the theoretical results.
\end{abstract}

\maketitle
\tableofcontents

% ============================================================
\section{Introduction}
% ============================================================

Differential geometry occupies a central place in modern mathematics, forming
the natural language for both pure mathematical structures and the fundamental
laws of physics. At its heart lies a deceptively simple question: how do we
measure and compare geometric quantities on curved spaces that cannot be
embedded in a flat ambient space in any canonical way?

The origins of the subject trace back to Carl Friedrich Gauss (1777--1855),
whose \emph{Disquisitiones generales circa superficies curvas} (1827)
introduced the revolutionary idea that the intrinsic geometry of a surface---
the geometry as experienced by a being living on the surface---is entirely
determined by the first fundamental form, without reference to any ambient
space. This is the content of Gauss's \emph{Theorema Egregium}: the Gaussian
curvature $K$ is an intrinsic invariant.

Bernhard Riemann (1826--1866), in his celebrated 1854 habilitation lecture
\emph{\"Uber die Hypothesen welche der Geometrie zu Grunde liegen}, generalised
Gauss's ideas to manifolds of arbitrary dimension. Riemann introduced what we
now call the Riemannian metric tensor $g_{ij}$, which allows the measurement
of lengths, angles, and volumes on smooth manifolds. His vision laid the
foundation for:
\begin{itemize}
  \item The intrinsic theory of smooth manifolds (Whitney, 1936),
  \item The Riemannian geometry of manifolds with a metric tensor,
  \item The covariant calculus of Gregorio Ricci-Curbastro and Tullio
        Levi-Civita (1900--1917),
  \item The general theory of relativity (Einstein, 1915), where spacetime
        is a 4-dimensional Lorentzian manifold.
\end{itemize}

The interplay between geometry and analysis is pervasive. Differential forms
connect curvature to topology via the de Rham cohomology. The spectrum of the
Laplace-Beltrami operator encodes geometric information. Minimal surfaces---
solving variational problems---arise both in soap-film experiments and in
modern string theory.

\subsection{Structure of This Paper}

Section~\ref{sec:manifolds} introduces smooth manifolds rigorously via charts
and atlases. Section~\ref{sec:metrics} develops Riemannian metrics, lengths,
and the exponential map. Section~\ref{sec:connections} treats connections and
covariant derivatives, culminating in the uniqueness of the Levi-Civita
connection. Section~\ref{sec:curvature} analyses the full hierarchy of
curvature tensors. Section~\ref{sec:gauss-bonnet} presents the Gauss-Bonnet
theorem and Chern's generalisation. Section~\ref{sec:geodesics} studies
geodesics, completeness, and Jacobi fields. Section~\ref{sec:submanifolds}
introduces the second fundamental form and minimal surfaces.
Section~\ref{sec:liegroups} connects Lie groups to geometry. Section~\ref{sec:applications}
describes applications to general relativity.

% ============================================================
\section{Smooth Manifolds}
\label{sec:manifolds}
% ============================================================

\begin{definition}[Topological Manifold]
A \emph{topological manifold} $M$ of dimension $n$ is a second-countable
Hausdorff topological space such that every point $p \in M$ has an open
neighbourhood $U$ homeomorphic to an open subset of $\mathbb{R}^n$.
\end{definition}

\begin{definition}[Chart and Atlas]
A \emph{chart} (or \emph{coordinate chart}) on $M$ is a pair $(U, \varphi)$
where $U \subseteq M$ is open and $\varphi: U \to V \subseteq \mathbb{R}^n$
is a homeomorphism. A \emph{smooth atlas} is a collection
$\mathcal{A} = \{(U_\alpha, \varphi_\alpha)\}$ of charts such that:
\begin{enumerate}[label=(\roman*)]
  \item $\bigcup_\alpha U_\alpha = M$ (covering condition),
  \item For any two charts $(U_\alpha, \varphi_\alpha)$ and
        $(U_\beta, \varphi_\beta)$ with $U_\alpha \cap U_\beta \neq \emptyset$,
        the \emph{transition map}
        \[
          \varphi_\beta \circ \varphi_\alpha^{-1}:
          \varphi_\alpha(U_\alpha \cap U_\beta) \to
          \varphi_\beta(U_\alpha \cap U_\beta)
        \]
        is a $C^\infty$ diffeomorphism.
\end{enumerate}
\end{definition}

\begin{definition}[Smooth Manifold]
A \emph{smooth manifold} is a topological manifold $M$ equipped with a maximal
smooth atlas $\mathcal{A}$ (also called a \emph{smooth structure}).
\end{definition}

\begin{example}
\begin{enumerate}[label=(\alph*)]
  \item The $n$-sphere $S^n = \{x \in \mathbb{R}^{n+1} : |x| = 1\}$ is a
        compact smooth $n$-manifold. The standard atlas consists of two
        stereographic projection charts.
  \item The torus $\mathbb{T}^2 = S^1 \times S^1$ is a compact 2-manifold
        with charts given by coordinate angle pairs $(\theta_1, \theta_2)$.
  \item The general linear group $\mathrm{GL}(n, \mathbb{R}) \subset M_n(\mathbb{R})$
        is an open submanifold of $\mathbb{R}^{n^2}$.
\end{enumerate}
\end{example}

\subsection{Tangent and Cotangent Spaces}

\begin{definition}[Tangent Space]
Let $p \in M$. A \emph{derivation} at $p$ is a linear map
$v: C^\infty(M) \to \mathbb{R}$ satisfying the Leibniz rule:
\[
  v(fg) = f(p)\,v(g) + g(p)\,v(f).
\]
The set of all derivations at $p$ is the \emph{tangent space} $T_pM$, which
is an $n$-dimensional real vector space. In a chart $(U, \varphi)$ with
coordinates $(x^1, \ldots, x^n)$, the \emph{coordinate basis vectors} are
\[
  \left.\frac{\partial}{\partial x^i}\right|_p f
  = \left.\frac{\partial (f \circ \varphi^{-1})}{\partial r^i}\right|_{\varphi(p)},
  \quad i = 1, \ldots, n.
\]
\end{definition}

\begin{definition}[Cotangent Space and Differential Forms]
The \emph{cotangent space} at $p$ is the dual $T_p^*M = (T_pM)^*$.
The differential of a smooth function $f: M \to \mathbb{R}$ is
$(df)_p \in T_p^*M$ defined by $(df)_p(v) = v(f)$ for $v \in T_pM$.
The coordinate differentials $\{dx^1|_p, \ldots, dx^n|_p\}$ form the
dual basis to $\{\partial/\partial x^i|_p\}$.
\end{definition}

\begin{definition}[Tangent Bundle]
The \emph{tangent bundle} is $TM = \bigsqcup_{p \in M} T_pM$ with the natural
projection $\pi: TM \to M$, $\pi(v) = p$ for $v \in T_pM$. A \emph{smooth
vector field} on $M$ is a smooth section $X: M \to TM$, i.e., $\pi \circ X = \mathrm{id}_M$.
\end{definition}

\subsection{Parametric Curves and the Frenet-Serret Frame}

A smooth curve in $\mathbb{R}^n$ is a smooth map $\mathbf{r}: I \to \mathbb{R}^n$.
For a regular curve ($|\mathbf{r}'(t)| > 0$), the \emph{Frenet-Serret frame}
$\{T, N, B\}$ consists of the unit tangent, principal normal, and binormal vectors.
The Frenet-Serret equations in $\mathbb{R}^3$ are:
\[
  T' = \kappa N, \quad N' = -\kappa T + \tau B, \quad B' = -\tau N,
\]
where $\kappa \geq 0$ is the \emph{curvature} and $\tau \in \mathbb{R}$ is
the \emph{torsion}. Explicitly:
\[
  \kappa = \frac{|\mathbf{r}' \times \mathbf{r}''|}{|\mathbf{r}'|^3},
  \qquad
  \tau = \frac{(\mathbf{r}' \times \mathbf{r}'') \cdot \mathbf{r}'''}{|\mathbf{r}' \times \mathbf{r}''|^2}.
\]

\begin{example}[Helix]
The helix $\mathbf{r}(t) = (\cos t, \sin t, t)$ has
$\kappa = 1/2$ and $\tau = 1/2$ (constant), illustrating a curve that
uniformly winds around a cylinder.
\end{example}

% ============================================================
\section{Riemannian Metrics}
\label{sec:metrics}
% ============================================================

\begin{definition}[Riemannian Metric]
A \emph{Riemannian metric} on a smooth manifold $M$ is a smooth assignment
of an inner product $g_p: T_pM \times T_pM \to \mathbb{R}$ to each tangent
space, varying smoothly with $p$. In local coordinates $(x^1, \ldots, x^n)$:
\[
  g = g_{ij}\,dx^i \otimes dx^j, \quad g_{ij} = g_{ji}, \quad (g_{ij}) \succ 0.
\]
The pair $(M, g)$ is called a \emph{Riemannian manifold}.
\end{definition}

\begin{example}[Standard metrics]
\begin{enumerate}[label=(\alph*)]
  \item \textbf{Euclidean space}: $g = \delta_{ij}\,dx^i\,dx^j = dx^2 + dy^2 + dz^2$.
  \item \textbf{Round sphere $S^2(r)$}: In spherical coordinates
        $(\theta, \phi)$, $g = r^2(d\theta^2 + \sin^2\!\theta\,d\phi^2)$.
  \item \textbf{Poincar\'e half-plane}: $\mathbb{H}^2 = \{(x,y): y>0\}$,
        $g = (dx^2 + dy^2)/y^2$. Constant curvature $K = -1$.
  \item \textbf{Torus $\mathbb{T}^2$}: Flat metric
        $g = d\theta_1^2 + d\theta_2^2$.
\end{enumerate}
\end{example}

\subsection{Arc Length and the First Fundamental Form}

For a parametric surface $\mathbf{r}(u,v): U \to \mathbb{R}^3$, the
\emph{first fundamental form} is
\[
  \mathrm{I} = E\,du^2 + 2F\,du\,dv + G\,dv^2,
\]
with coefficients $E = \mathbf{r}_u \cdot \mathbf{r}_u$,
$F = \mathbf{r}_u \cdot \mathbf{r}_v$, $G = \mathbf{r}_v \cdot \mathbf{r}_v$.
The \emph{arc length} of a curve $\mathbf{r}(t)$ is:
\[
  L = \int_{t_0}^{t_1} |\mathbf{r}'(t)|\,dt
    = \int_{t_0}^{t_1} \sqrt{g_{ij}\,\dot{x}^i\,\dot{x}^j}\,dt.
\]

\subsection{The Exponential Map}

\begin{definition}[Exponential Map]
Let $(M, g)$ be a Riemannian manifold and $p \in M$. For $v \in T_pM$,
let $\gamma_v$ denote the unique geodesic with $\gamma_v(0) = p$ and
$\gamma_v'(0) = v$. The \emph{exponential map} at $p$ is
\[
  \exp_p: T_pM \supseteq \mathcal{D}_p \to M, \quad
  \exp_p(v) = \gamma_v(1),
\]
defined on the maximal domain $\mathcal{D}_p$ where $\gamma_v$ is defined
on $[0,1]$.
\end{definition}

\begin{proposition}
For any $p \in M$, the exponential map $\exp_p$ is a local diffeomorphism
from a neighbourhood of $0 \in T_pM$ to a neighbourhood of $p \in M$.
\end{proposition}

Normal coordinates centred at $p$, obtained via $\exp_p$, satisfy
$g_{ij}(p) = \delta_{ij}$ and $\partial_k g_{ij}(p) = 0$, which is
fundamental for many local computations.

% ============================================================
\section{Connections and Covariant Derivatives}
\label{sec:connections}
% ============================================================

A connection provides the means to differentiate vector fields along curves
on a manifold, generalising the directional derivative of vector fields
in $\mathbb{R}^n$.

\begin{definition}[Affine Connection]
An \emph{affine connection} on $M$ is a map
$\nabla: \mathfrak{X}(M) \times \mathfrak{X}(M) \to \mathfrak{X}(M)$,
$(X, Y) \mapsto \nabla_X Y$, satisfying:
\begin{enumerate}[label=(\roman*)]
  \item $\nabla_{fX+gY} Z = f \nabla_X Z + g \nabla_Y Z$ \quad ($C^\infty$-linearity in $X$),
  \item $\nabla_X(Y + Z) = \nabla_X Y + \nabla_X Z$,
  \item $\nabla_X(fY) = (Xf)\,Y + f\,\nabla_X Y$ (Leibniz rule).
\end{enumerate}
In local coordinates, $\nabla_{\partial_i} \partial_j = \Gamma^k_{ij} \partial_k$,
where $\Gamma^k_{ij}$ are the \emph{Christoffel symbols}.
\end{definition}

\begin{definition}[Metric Compatibility and Torsion-Free]
A connection $\nabla$ on $(M, g)$ is:
\begin{itemize}
  \item \emph{metric compatible} if $\nabla g = 0$, i.e.,
        $X\bigl(g(Y,Z)\bigr) = g(\nabla_X Y, Z) + g(Y, \nabla_X Z)$,
  \item \emph{torsion-free} (symmetric) if $T(X,Y) = \nabla_X Y - \nabla_Y X - [X,Y] = 0$.
\end{itemize}
\end{definition}

\begin{theorem}[Levi-Civita Connection --- Fundamental Theorem of Riemannian Geometry]
\label{thm:levi-civita}
Let $(M, g)$ be a Riemannian manifold. There exists a \textbf{unique} affine
connection $\nabla$ on $M$ that is both metric compatible and torsion-free.
This connection is called the \emph{Levi-Civita connection}.
\end{theorem}

\begin{proof}
\textbf{Uniqueness.} Suppose $\nabla$ is metric compatible and torsion-free.
Using metric compatibility cyclically:
\begin{align*}
  Xg(Y,Z) &= g(\nabla_X Y, Z) + g(Y, \nabla_X Z), \\
  Yg(Z,X) &= g(\nabla_Y Z, X) + g(Z, \nabla_Y X), \\
  Zg(X,Y) &= g(\nabla_Z X, Y) + g(X, \nabla_Z Y).
\end{align*}
Adding the first two and subtracting the third, then using torsion-freeness
$\nabla_X Y - \nabla_Y X = [X,Y]$, yields the \emph{Koszul formula}:
\begin{multline*}
  2\,g(\nabla_X Y, Z) = Xg(Y,Z) + Yg(X,Z) - Zg(X,Y) \\
    + g([X,Y], Z) - g([X,Z], Y) - g([Y,Z], X).
\end{multline*}
Since $g$ is non-degenerate, $\nabla_X Y$ is uniquely determined. This proves
uniqueness.

\textbf{Existence.} Define $\nabla_X Y$ by the Koszul formula. One verifies
directly that this satisfies all axioms of a connection, is metric compatible,
and torsion-free. The corresponding Christoffel symbols are
\[
  \Gamma^k_{ij} = \tfrac{1}{2} g^{kl}
  \bigl(\partial_i g_{jl} + \partial_j g_{il} - \partial_l g_{ij}\bigr).
\]
\end{proof}

\subsection{Parallel Transport}

\begin{definition}[Parallel Transport]
A vector field $V$ along a curve $\gamma: [a,b] \to M$ is
\emph{parallel} (along $\gamma$) if $\nabla_{\gamma'}V = 0$.
In coordinates:
\[
  \frac{dV^k}{dt} + \Gamma^k_{ij}\,\frac{dx^i}{dt}\,V^j = 0.
\]
\emph{Parallel transport} along $\gamma$ is the isomorphism
$P_{a \to t}: T_{\gamma(a)}M \to T_{\gamma(t)}M$ mapping each tangent
vector at $\gamma(a)$ to its parallel translate at $\gamma(t)$.
\end{definition}

\begin{remark}
Parallel transport on a curved manifold is path-dependent. The holonomy of a
closed loop on $S^2$ gives a rotation equal to the solid angle enclosed,
directly visible as the Foucault pendulum precession on Earth.
\end{remark}

% ============================================================
\section{Curvature}
\label{sec:curvature}
% ============================================================

Curvature measures the failure of second-order derivatives to commute,
which is the infinitesimal version of path-dependent parallel transport.

\begin{definition}[Riemann Curvature Tensor]
The \emph{Riemann curvature tensor} is the $(1,3)$-tensor:
\[
  R(X,Y)Z = \nabla_X \nabla_Y Z - \nabla_Y \nabla_X Z - \nabla_{[X,Y]} Z.
\]
In local coordinates, the components $R^l{}_{\!kij}$ are defined by
$R(\partial_i, \partial_j)\partial_k = R^l{}_{\!kij}\,\partial_l$, giving:
\[
  R^l{}_{\!kij} = \partial_i \Gamma^l_{jk} - \partial_j \Gamma^l_{ik}
  + \Gamma^l_{im}\Gamma^m_{jk} - \Gamma^l_{jm}\Gamma^m_{ik}.
\]
\end{definition}

\begin{proposition}[Symmetries of the Riemann Tensor]
With indices lowered $R_{ijkl} = g_{im}R^m{}_{\!jkl}$:
\begin{enumerate}[label=(\roman*)]
  \item $R_{ijkl} = -R_{jikl}$ (antisymmetry in first pair),
  \item $R_{ijkl} = -R_{ijlk}$ (antisymmetry in second pair),
  \item $R_{ijkl} = R_{klij}$ (pair symmetry),
  \item $R_{ijkl} + R_{iklj} + R_{iljk} = 0$ (algebraic Bianchi identity).
\end{enumerate}
\end{proposition}

\subsection{Ricci and Scalar Curvature}

\begin{definition}
The \emph{Ricci tensor} $\mathrm{Ric}$ is the contraction:
\[
  R_{ij} = R^k{}_{\!ikj} = g^{kl}R_{kilj}.
\]
The \emph{scalar curvature} is the full trace:
\[
  R = g^{ij}R_{ij} = \sum_i R^i{}_i.
\]
\end{definition}

\begin{definition}[Sectional Curvature]
For a 2-plane $\sigma = \mathrm{span}\{u, v\} \subset T_pM$, the
\emph{sectional curvature} is:
\[
  K(\sigma) = \frac{g(R(u,v)v, u)}{g(u,u)g(v,v) - g(u,v)^2}.
\]
\end{definition}

\begin{example}[Curvature of classical surfaces]
\begin{enumerate}[label=(\alph*)]
  \item \textbf{Sphere $S^n(r)$}: Constant sectional curvature $K = 1/r^2$.
  \item \textbf{Euclidean space}: $K = 0$ identically.
  \item \textbf{Hyperbolic space $\mathbb{H}^n$}: Constant $K = -1$.
  \item \textbf{Torus $\mathbb{T}^2$} (flat metric): $K = 0$.
\end{enumerate}
\end{example}

\subsection{Gaussian Curvature via Fundamental Forms}

For a surface $\mathbf{r}(u,v)$ in $\mathbb{R}^3$, the \emph{second
fundamental form} has coefficients:
\[
  L = \mathbf{r}_{uu} \cdot \mathbf{n}, \quad
  M = \mathbf{r}_{uv} \cdot \mathbf{n}, \quad
  N = \mathbf{r}_{vv} \cdot \mathbf{n},
\]
where $\mathbf{n} = (\mathbf{r}_u \times \mathbf{r}_v)/|\mathbf{r}_u \times \mathbf{r}_v|$
is the unit normal. The \emph{Gaussian curvature} and \emph{mean curvature} are:
\[
  K = \frac{LN - M^2}{EG - F^2}, \qquad
  H = \frac{EN + GL - 2FM}{2(EG - F^2)}.
\]

\begin{theorem}[Theorema Egregium --- Gauss, 1827]
The Gaussian curvature $K$ of a surface in $\mathbb{R}^3$ depends only on
the first fundamental form $E, F, G$ and their derivatives. In particular,
$K$ is invariant under local isometries.
\end{theorem}

\begin{remark}
This theorem is genuinely remarkable (as its Latin name suggests): a flat
piece of paper rolled into a cylinder has $K = 0$ both flat and cylindrical,
while $K \neq 0$ for a sphere cannot be developed (mapped isometrically)
onto a flat surface, explaining why flat maps of the Earth necessarily
distort areas or angles.
\end{remark}

% ============================================================
\section{The Gauss-Bonnet Theorem}
\label{sec:gauss-bonnet}
% ============================================================

The Gauss-Bonnet theorem is one of the most beautiful results in mathematics,
connecting the local differential-geometric concept of curvature to the global
topological invariant---the Euler characteristic.

\begin{theorem}[Gauss-Bonnet for Closed Surfaces]
\label{thm:gauss-bonnet}
Let $(M, g)$ be a compact oriented Riemannian 2-manifold without boundary.
Then:
\[
  \int_M K\,dA = 2\pi\,\chi(M),
\]
where $K$ is the Gaussian curvature, $dA$ the area element, and $\chi(M)$
the Euler characteristic of $M$.
\end{theorem}

\begin{proof}[Proof sketch]
We proceed by a combinatorial-analytic argument.

\textbf{Step 1 (Local Gauss-Bonnet).} For a geodesic polygon $P$ with
vertices $p_1, \ldots, p_k$, interior angles $\alpha_i$, and exterior angles
$\theta_i = \pi - \alpha_i$:
\[
  \int_P K\,dA + \oint_{\partial P} \kappa_g\,ds = \sum_{i=1}^k (\pi - \alpha_i),
\]
where $\kappa_g$ is the geodesic curvature of $\partial P$.
This follows from the Gauss map and Stokes' theorem.

\textbf{Step 2 (Triangulation).} Triangulate $M$ by geodesic triangles
$T_1, \ldots, T_F$ with $V$ vertices and $E$ edges. Since each triangle has
3 edges and each internal edge is shared by 2 triangles: $3F = 2E$ for an
interior triangulation. Apply the local formula to each triangle.

\textbf{Step 3 (Summation).} Summing over all triangles:
\[
  \int_M K\,dA = \sum_{i} \int_{T_i} K\,dA
  = (V - E + F) \cdot 2\pi = 2\pi\,\chi(M),
\]
using the Euler formula $\chi(M) = V - E + F$ and the fact that
interior angles around each vertex sum to $2\pi$.
\end{proof}

\begin{corollary}[Topological consequences]
\begin{enumerate}[label=(\roman*)]
  \item For $S^2$: $\chi(S^2) = 2$, so $\int_{S^2} K\,dA = 4\pi = 4\pi r^2 \cdot \frac{1}{r^2}$.
  \item A torus $\mathbb{T}^2$ has $\chi = 0$, so $\int K\,dA = 0$: the
        positive and negative curvature contributions exactly cancel.
  \item No metric on $S^2$ can have $K \leq 0$ everywhere (with strict
        inequality somewhere), and no metric on $\mathbb{T}^2$ can have $K > 0$.
\end{enumerate}
\end{corollary}

\subsection{Chern-Gauss-Bonnet Theorem}

\begin{theorem}[Chern-Gauss-Bonnet]
\label{thm:chern-gauss-bonnet}
Let $M$ be a compact oriented Riemannian manifold of even dimension
$n = 2m$ without boundary. Then:
\[
  \chi(M) = \frac{1}{(2\pi)^m \, 2^m \, m!} \int_M \mathrm{Pf}(\Omega),
\]
where $\Omega$ is the curvature 2-form of the Levi-Civita connection
and $\mathrm{Pf}$ denotes the Pfaffian of an antisymmetric matrix.
For $n=4$ with Weyl tensor $W$ and Ricci tensor $\mathrm{Ric}$:
\[
  \chi(M) = \frac{1}{8\pi^2}\int_M
  \left(|W|^2 - \frac{1}{2}|\overset{\circ}{\mathrm{Ric}}|^2
  + \frac{R^2}{24}\right) dV,
\]
where $\overset{\circ}{\mathrm{Ric}}$ is the traceless part of the Ricci tensor.
\end{theorem}

This was proved by Shiing-Shen Chern in 1944 using the formalism of
characteristic classes and the Gauss-Bonnet-Chern form, without requiring
any embedding into an ambient Euclidean space.

% ============================================================
\section{Geodesics}
\label{sec:geodesics}
% ============================================================

Geodesics are the ``straight lines'' of Riemannian geometry---curves that
locally minimise length.

\begin{definition}[Geodesic]
A smooth curve $\gamma: I \to M$ is a \emph{geodesic} if its tangent vector
is parallel along itself:
\[
  \nabla_{\gamma'}\gamma' = 0.
\]
In local coordinates, with $\gamma(t) = (x^1(t), \ldots, x^n(t))$, this
is the \emph{geodesic equation}:
\[
  \ddot{x}^k + \Gamma^k_{ij}\,\dot{x}^i\,\dot{x}^j = 0,
  \quad k = 1, \ldots, n.
\]
\end{definition}

\begin{proposition}
Geodesics are constant-speed curves: $|\gamma'(t)|_g = \text{const}$.
\end{proposition}

\begin{example}[Geodesics on classical spaces]
\begin{enumerate}[label=(\alph*)]
  \item In $\mathbb{R}^n$: straight lines $\gamma(t) = p + tv$.
  \item On $S^2$: great circles (intersections of planes through the origin).
  \item In $\mathbb{H}^2$: vertical lines $x = \text{const}$ and semicircles
        with diameter on the $x$-axis.
\end{enumerate}
\end{example}

\subsection{Hopf-Rinow Theorem}

\begin{theorem}[Hopf-Rinow]
\label{thm:hopf-rinow}
Let $(M, g)$ be a connected Riemannian manifold. The following are equivalent:
\begin{enumerate}[label=(\roman*)]
  \item $(M, d_g)$ is a complete metric space (every Cauchy sequence converges),
  \item $M$ is \emph{geodesically complete}: for every $p \in M$ and every
        $v \in T_pM$, the geodesic $\gamma_v$ is defined for all $t \in \mathbb{R}$,
  \item For some (equivalently, every) point $p \in M$, the exponential map
        $\exp_p$ is defined on all of $T_pM$.
\end{enumerate}
Moreover, any of these conditions implies:
\begin{enumerate}[label=(\alph*)]
  \item Any two points in $M$ are connected by a minimising geodesic.
\end{enumerate}
\end{theorem}

\begin{proof}[Proof sketch]
The implication (i) $\Rightarrow$ (ii) uses the fact that geodesics cannot
``reach infinity in finite time'' if the metric space is complete: if a
geodesic $\gamma$ were defined only on $[0, T)$, then $\gamma(t_n)$ for
$t_n \nearrow T$ would be Cauchy, hence convergent to some $p^* \in M$,
and the geodesic extends through $p^*$.

The implication (iii) $\Rightarrow$ (a) is proved by showing that for any
$q \in M$, the infimum of the length functional over all curves from $p$
to $q$ is attained: one constructs a minimising sequence of curves and
extracts a convergent subsequence using the geodesic ball structure, then
uses $\exp_p$ defined on all of $T_pM$ to extend the minimising geodesic.
\end{proof}

\subsection{Jacobi Fields}

\begin{definition}[Jacobi Field]
Let $\gamma$ be a geodesic. A vector field $J$ along $\gamma$ is a
\emph{Jacobi field} if it satisfies the \emph{Jacobi equation}:
\[
  \nabla^2_{\gamma'} J + R(\gamma', J)\gamma' = 0,
\]
i.e., $J'' + R(\gamma', J)\gamma' = 0$ where $J'' = \nabla_{\gamma'}\nabla_{\gamma'} J$.
\end{definition}

Jacobi fields describe infinitesimal geodesic variations and encode how
nearby geodesics converge or diverge. A \emph{conjugate point} to $p$
along $\gamma$ is a point $\gamma(t_0)$ where a non-trivial Jacobi field
vanishes. On $S^2$, the conjugate locus of the north pole is the south pole.

% ============================================================
\section{Submanifolds}
\label{sec:submanifolds}
% ============================================================

\begin{definition}[Embedded Submanifold]
An \emph{embedded submanifold} of $(M, g)$ is a smooth submanifold
$\iota: \Sigma \hookrightarrow M$ with the induced (pullback) metric
$\bar{g} = \iota^* g$.
\end{definition}

\subsection{The Second Fundamental Form}

For a hypersurface $\Sigma \subset M$ with unit normal $\mathbf{n}$,
the \emph{second fundamental form} is the symmetric bilinear form:
\[
  \mathrm{II}(X,Y) = g(\nabla_X Y, \mathbf{n}), \quad X, Y \in T\Sigma,
\]
encoding how $\Sigma$ bends in $M$. Its trace (with respect to $\bar{g}$)
is the \emph{mean curvature}:
\[
  H = \frac{1}{n-1}\,\mathrm{tr}_{\bar{g}}(\mathrm{II}).
\]

\begin{definition}[Minimal Surface]
A surface $\Sigma$ is \emph{minimal} if $H \equiv 0$.
\end{definition}

\begin{example}[Classical minimal surfaces]
\begin{enumerate}[label=(\alph*)]
  \item \textbf{Catenoid}: $\mathbf{r}(u,v) = (\cosh u \cos v, \cosh u \sin v, u)$.
        The unique minimal surface of revolution.
  \item \textbf{Helicoid}: $\mathbf{r}(u,v) = (u\cos v, u\sin v, v)$.
        Simply connected minimal surface, isometric to the catenoid.
  \item \textbf{Enneper surface}: $\mathbf{r}(u,v)$ given by Weierstrass-Enneper
        representation; self-intersecting but minimal.
\end{enumerate}
\end{example}

\subsection{The Gauss and Codazzi Equations}

The Gauss equation relates the intrinsic curvature of $\Sigma$ to the
ambient curvature and the second fundamental form:
\[
  \bar{R}(X,Y,Z,W) = R(X,Y,Z,W)
  + \mathrm{II}(X,W)\,\mathrm{II}(Y,Z) - \mathrm{II}(X,Z)\,\mathrm{II}(Y,W).
\]
In the classical case $M = \mathbb{R}^3$, $R = 0$, this gives
$K = \kappa_1 \kappa_2$ (product of principal curvatures), which is the
content of Theorema Egregium.

The Codazzi-Mainardi equations impose integrability conditions:
\[
  (\nabla_X \mathrm{II})(Y,Z) = (\nabla_Y \mathrm{II})(X,Z) + R(X,Y)Z^{\perp}.
\]

% ============================================================
\section{Lie Groups and Geometry}
\label{sec:liegroups}
% ============================================================

Lie groups combine the algebraic structure of a group with the smooth
structure of a manifold, providing a rich source of symmetric spaces.

\begin{definition}[Lie Group]
A \emph{Lie group} is a group $G$ that is simultaneously a smooth manifold,
such that the group operations $\mu: G \times G \to G$, $(g,h) \mapsto gh$
and $\iota: G \to G$, $g \mapsto g^{-1}$ are smooth.
\end{definition}

\begin{example}
\begin{enumerate}[label=(\alph*)]
  \item $\mathbb{R}^n$ (addition), $\mathbb{R}^* = \mathbb{R} \setminus \{0\}$,
  \item $\mathrm{GL}(n,\mathbb{R})$, $\mathrm{SL}(n,\mathbb{R})$, $\mathrm{O}(n)$, $\mathrm{SO}(n)$, $\mathrm{U}(n)$, $\mathrm{SU}(n)$,
  \item The circle group $S^1 \cong \mathrm{U}(1)$.
\end{enumerate}
\end{example}

\begin{definition}[Lie Algebra]
The \emph{Lie algebra} $\mathfrak{g}$ of $G$ is $T_eG$ equipped with the
bracket $[X, Y] = \mathcal{L}_X Y$ (Lie derivative). For matrix Lie groups,
$\mathfrak{g}$ consists of matrices and $[A,B] = AB - BA$.
\end{definition}

\subsection{Left-Invariant Metrics and Killing Fields}

A Riemannian metric $g$ on $G$ is \emph{left-invariant} if
$L_h^* g = g$ for all $h \in G$, where $L_h(x) = hx$. Such a metric is
uniquely determined by its value at $e$, i.e., by an inner product on $\mathfrak{g}$.

\begin{definition}[Killing Vector Field]
A vector field $X$ on $(M, g)$ is a \emph{Killing field} if
$\mathcal{L}_X g = 0$, i.e., $X$ generates isometries. Equivalently:
\[
  \nabla_i X_j + \nabla_j X_i = 0.
\]
\end{definition}

\begin{proposition}
On a Lie group $G$ with left-invariant metric, the left-invariant vector
fields are Killing fields if and only if the metric is also right-invariant
(\emph{bi-invariant}).
\end{proposition}

\begin{example}
The bi-invariant metric on $\mathrm{SO}(3)$ is $\langle A, B \rangle = -\tfrac{1}{2}\mathrm{tr}(AB)$.
Geodesics are one-parameter subgroups $t \mapsto e^{tX}$, and sectional
curvature satisfies $K(\sigma) \geq 0$.
\end{example}

% ============================================================
\section{Applications to General Relativity}
\label{sec:applications}
% ============================================================

Albert Einstein's general theory of relativity (1915) is arguably the
greatest application of differential geometry to physics. Spacetime is
modelled as a 4-dimensional \emph{Lorentzian manifold} $(M, g)$ with
metric signature $(-,+,+,+)$, and gravity is identified with spacetime
curvature.

\subsection{Lorentzian Manifolds}

\begin{definition}
A \emph{Lorentzian manifold} $(M, g)$ is a smooth manifold with a metric
tensor $g$ of signature $(-,+,+,+)$. Vectors $v \in T_pM$ are classified as:
\begin{itemize}
  \item \emph{timelike}: $g(v,v) < 0$ (massive particles),
  \item \emph{null (lightlike)}: $g(v,v) = 0$ (photons),
  \item \emph{spacelike}: $g(v,v) > 0$.
\end{itemize}
\end{definition}

\subsection{Einstein Field Equations}

\begin{definition}[Einstein Tensor]
The \emph{Einstein tensor} is:
\[
  G_{\mu\nu} = R_{\mu\nu} - \tfrac{1}{2} R\,g_{\mu\nu},
\]
where $R_{\mu\nu}$ is the Ricci tensor and $R$ the scalar curvature.
It satisfies the contracted Bianchi identity $\nabla^\mu G_{\mu\nu} = 0$.
\end{definition}

\begin{theorem}[Einstein Field Equations]
In general relativity, the geometry of spacetime is determined by the
distribution of matter and energy through the equations:
\[
  G_{\mu\nu} + \Lambda g_{\mu\nu} = \frac{8\pi G}{c^4}\,T_{\mu\nu},
\]
where $\Lambda$ is the cosmological constant, $G$ Newton's gravitational
constant, and $T_{\mu\nu}$ the stress-energy tensor. In vacuum ($T_{\mu\nu} = 0$,
$\Lambda = 0$), this reduces to $R_{\mu\nu} = 0$ (\emph{Ricci-flat}).
\end{theorem}

\subsection{Geodesics as Worldlines}

\begin{proposition}
In general relativity, freely falling particles (not subject to any non-gravitational
forces) move along \emph{timelike geodesics} of the spacetime metric. Photons
travel along \emph{null geodesics}.
\end{proposition}

\begin{example}[Schwarzschild Geodesics]
The Schwarzschild metric (exterior to a spherically symmetric mass $m$) is:
\[
  ds^2 = -\left(1 - \frac{r_s}{r}\right)c^2\,dt^2
  + \left(1 - \frac{r_s}{r}\right)^{-1}dr^2
  + r^2(d\theta^2 + \sin^2\!\theta\,d\phi^2),
\]
with Schwarzschild radius $r_s = 2Gm/c^2$. Circular geodesics exist at
$r = 3r_s/2$ (photon sphere) and $r = 3r_s$ (innermost stable circular orbit).
The precession of Mercury's perihelion ($43''$ per century) is a direct
prediction from the geodesic equations.
\end{example}

\subsection{Connection to the Python Module}

The accompanying module \texttt{differential\_geometry.py} implements:
\begin{itemize}
  \item \texttt{ParametricCurve}: Frenet-Serret frame, curvature $\kappa$, torsion $\tau$,
  \item \texttt{ParametricSurface}: first and second fundamental forms, $K$, $H$,
  \item \texttt{RiemannianGeometry}: Christoffel symbols $\Gamma^k_{ij}$,
        Riemann tensor $R^l{}_{kij}$, Ricci tensor $R_{ij}$, scalar curvature,
  \item \texttt{GeodesicComputation}: numerical geodesic integration via RK4/RK45,
  \item \texttt{ClassicalSurfaces}: sphere, torus, catenoid, helicoid, saddle.
\end{itemize}
Curvature computations use the Koszul formula for Christoffel symbols and
contract accordingly. Geodesics are solved as second-order ODEs using
\texttt{scipy.integrate.solve\_ivp} with the Dormand-Prince method.

% ============================================================
\section*{References}
% ============================================================

\begin{thebibliography}{99}

\bibitem{docarmo1976}
M.~P. do~Carmo,
\textit{Differential Geometry of Curves and Surfaces},
Prentice-Hall, 1976.

\bibitem{docarmo1992}
M.~P. do~Carmo,
\textit{Riemannian Geometry},
Birkh\"auser, Boston, 1992.

\bibitem{spivak1970}
M.~Spivak,
\textit{A Comprehensive Introduction to Differential Geometry},
Vol.~1--5, Publish or Perish, 1970--1975.

\bibitem{lee1997}
J.~M. Lee,
\textit{Riemannian Manifolds: An Introduction to Curvature},
Springer, New York, 1997.

\bibitem{lee2013}
J.~M. Lee,
\textit{Introduction to Smooth Manifolds},
2nd ed., Springer, New York, 2013.

\bibitem{oneill1983}
B.~O'Neill,
\textit{Semi-Riemannian Geometry with Applications to Relativity},
Academic Press, 1983.

\bibitem{petersen2006}
P.~Petersen,
\textit{Riemannian Geometry},
2nd ed., Springer, New York, 2006.

\bibitem{chern1944}
S.-S. Chern,
A simple intrinsic proof of the Gauss-Bonnet formula for closed Riemannian manifolds,
\textit{Ann. of Math.}, \textbf{45} (1944), 747--752.

\bibitem{gallot1990}
S.~Gallot, D.~Hulin, J.~Lafontaine,
\textit{Riemannian Geometry},
3rd ed., Springer, Berlin, 2004.

\bibitem{hopfrinow1931}
H.~Hopf, W.~Rinow,
\"Uber den Begriff der vollst\"andigen differentialgeometrischen Fl\"ache,
\textit{Comment. Math. Helv.}, \textbf{3} (1931), 209--225.

\end{thebibliography}

\end{document}
